\documentclass[letterpaper,single,8pt]{article}
\usepackage[margin=2.2cm]{geometry}
\usepackage{xspace}
\usepackage[hidelinks, plainpages=false, pdfpagelabels]{hyperref}

\newcommand{\chap}{Chapter\xspace}
\newcommand{\ie}{i.e.,\xspace}
\newcommand{\eg}{e.g.,\xspace}
\newcommand{\restler}{RESTler\xspace}
\newcommand{\pythia}{Pythia\xspace}
\def\gitlab{GitLab\xspace}


\begin{document}
\title{\vspace{-1.2cm}{\bf Testing of Cloud Services}}
\author{ {\rm Vaggelis Atlidakis} }
\date{}
\maketitle

Over the last decade, we have witnessed an explosion in cloud services
for hosting software applications (Software-as-a-Service), for
building distributed services (Platform-as-a-Service),
and for providing general computing infrastructure
(Infrastructure-as-a-Service). Thousands of new cloud services have been
deployed by cloud platform providers, such as Amazon Web Services~\cite{aws},
Google Cloud~\cite{google-cloud}, and Microsoft Azure~\cite{azure}, and support
customers who are modernizing their processes by switching from the complexity
of owning and maintaining their own, on-premise Information Technology (IT)
infrastructure to instead
simply access and pay on demand cutting edge technologies. According to a recent
study, in $2020$, the public cloud services market is expected to reach around
$266$ billion U.S. dollars in size and by $2022$, its market revenue is forecast
to exceed $350$ billion U.S. dollars~\cite{cloudmarket}.

Today, most cloud services are programmatically accessed by third-party
applications~\cite{REST-book} and other services~\cite{services-book}
through Application Programming Interfaces (APIs) that follow the
REpresentational State Transfer (REST) software architectural style~\cite{REST}.
REST APIs are implemented
on top of the ubiquitous HTTP and HTTPS protocols, and provide requesting
systems (clients) with a uniform and predefined set of stateless operations
in order to create, monitor, manage, and delete cloud resources.
Furthermore, cloud service developers use interface-description languages,
such as the OpenAPI Specification (OAS)~\cite{openapi}, to describe and document
their REST APIs, including what requests the service can handle through its REST
API and the format of those requests, what responses may be received, and the
respective response format~\cite{azureapis,googleapis,awsapis}.

The fact that the vast majority of cloud services are accessed through REST APIs
that are well-documented with interface-specification languages
presents a unique opportunity to build systems that automatically
test cloud services through their APIs. When is a cloud service reliable
and secure to be publicly deployed? What kinds of software errors may be hiding
behind REST APIs? And how critical these errors may be?
Automatically answering such questions for production-scale, cloud services
still remains an open research challenge and is of paramount importance for a
multibillion-dollar market where competing cloud providers seek to avoid
reliability and security incidents that will attract negative publicity.

Indeed, despite the rapidly evolving cloud ecosystem, systems for automatically
testing cloud services through their REST APIs are still in their
infancy. The most sophisticated testing systems currently available for REST APIs
capture live API traffic, and then fuzz and replay the recorded traffic with
the hope of finding software errors. Many of these fuzzing systems were born as
extensions of website testing and scanning
tools~\cite{appspider,qualysWAS,apifuzzer}.
Since these REST API testing systems are all recent and not yet widely used, it is
still largely unclear how effective they are in finding unknown errors,
especially those that hide behind complex APIs and require sequences of events
in order to be exposed.

My work leverages the structured usage of cloud services through REST APIs
and feedback obtained during interaction with cloud services
in the form of response status codes and response objects
to build systems that test cloud services through their APIs in
an {\em automatic}, {\em efficient}, and {\em learning-based way}
~\cite{restler,restler2,pythia,atlidakis}.

Along with my collaborators, I have built systems that are
automatic in that they require minimal manual intervention to test
target cloud services, efficient in that they find within a reasonable
time-frame (e.g., in few hours) previously-unknown software errors that were
beyond the reach of past tools, and learning-based in that they learn without
predefined rules or heuristics how to test target cloud services.
Test automation is achieved by leveraging the common structure present in the
ecosystem of cloud services that are accessed through well-documented REST
APIs. Efficiency is achieved by careful consideration of the semantics of the
RESTful architectural design style, which allows to generate test sequences that
consist of multiple API requests and test target services deeply, and by
utilizing feedback obtained during interaction with the target services in order
to prune large search spaces and avoid redundant test cases. Finally, the
systems I build pursue the avenue of learning-based program analysis and leverage
recent advancements in the area of machine learning~\cite{seq2seq1,seq2seq2}
in order to learn from past tests common usage patterns of target cloud services
and generate new tests.

My systems have been used to test
various production-scale, open-source and proprietary cloud services ,such as \gitlab~\cite{gitlab},
and Azure~\cite{azure} and Office-365~\cite{office365} cloud services,
and have lead to tens of bug fixes so far.

\pagebreak
\bibliographystyle{abbrv}
\bibliography{statement}


\end{document}
